\documentclass[conference]{IEEEtran}
\usepackage{amsmath,amssymb}
\usepackage{graphicx}
\usepackage{booktabs}
\usepackage{cite}

\begin{document}

\title{Paper Title: Concise and Informative}

\author{
  \IEEEauthorblockN{First Author}
  \IEEEauthorblockA{Department, University\\
  City, Country\\
  first.author@example.edu}
  \and
  \IEEEauthorblockN{Second Author}
  \IEEEauthorblockA{Department, University\\
  City, Country\\
  second.author@example.edu}
}

\maketitle

\begin{abstract}
Summarise the problem, your approach, the main result and why it matters,
in about 150 to 250 words.
\end{abstract}

\begin{IEEEkeywords}
first keyword, second keyword, third keyword
\end{IEEEkeywords}

\section{Introduction}
State the problem and its context. Explain what is missing in prior work and
list your contributions~\cite{example2024}.

\section{Related Work}
Group earlier work by theme and say how yours differs.

\section{Method}
Describe your approach. Number equations you refer to later:
\begin{equation}
  \mathcal{L}(\theta) = \frac{1}{N}\sum_{i=1}^{N} \ell\bigl(f_\theta(x_i), y_i\bigr).
  \label{eq:loss}
\end{equation}

\section{Results}
\begin{table}[t]
  \centering
  \caption{Main results.}
  \label{tab:results}
  \begin{tabular}{lcc}
    \toprule
    Method & Accuracy (\%) & Latency (ms) \\
    \midrule
    Baseline & 91.2 & 34 \\
    Proposed & 93.8 & 19 \\
    \bottomrule
  \end{tabular}
\end{table}
Table~\ref{tab:results} compares the proposed method with the baseline.

\section{Conclusion}
Summarise the findings and suggest future work.

\bibliographystyle{IEEEtran}
\bibliography{references}

\end{document}
